\chapter{Large deviations}
\label{chap:prereq_ldp}

Library-shaped large deviation theory, for \texttt{ForLMLExplore/Probability/LargeDeviations/}.
It is the subject of the LMLRoadmaps target Identification, Layer 8.3; no large deviation
principle exists in Mathlib, Tau Ceti, LML or LMLExplore.

\begin{definition}
  \label{def:good_rate}
  \lean{LMLExplore.IsGoodRateFunction}
  \leanok
  A good rate function $I : X \to [0, \infty]$ is lower semicontinuous with compact sublevel sets
  $\{x \mid I(x) \le c\}$, $c < \infty$ \cite{DZ98}.
\end{definition}

\begin{definition}
  \label{def:ldp}
  \lean{LMLExplore.IsLDP}
  \leanok
  A sequence of measures $(P_n)$ on a topological space $X$ satisfies the large deviation
  principle with speed $n$ and rate function $I$ if $I$ is lower semicontinuous and for every set
  $S$,
  $-\inf_{\inter S} I \le \liminf_n \frac1n \log P_n(S)$ and
  $\limsup_n \frac1n \log P_n(S) \le -\inf_{\cl S} I$ \cite{DZ98}.
\end{definition}

\begin{lemma}
  \label{lem:ldp_of_bounds}
  \lean{LMLExplore.IsLDP.of_isClosed_of_mem_nhds,
    LMLExplore.limsup_log_div_le_neg_of_le_mul_exp,
    LMLExplore.neg_le_liminf_log_div_of_mul_exp_le}
  \leanok
  \uses{def:ldp}
  If $I$ is lower semicontinuous, $\limsup_n \frac1n \log P_n(F) \le -\inf_F I$ for every closed
  set $F$, and $-I(x) \le \liminf_n \frac1n \log P_n(U)$ for every neighbourhood $U$ of every point
  $x$, then $(P_n)$ satisfies the large deviation principle with rate function $I$. Bounds
  $u_n \le C e^{-n r}$ for every $r < c$ give $\limsup_n \frac1n \log u_n \le -c$, and bounds
  $C e^{-n r} \le u_n$ with $C > 0$ for every $r > c$ give $-c \le \liminf_n \frac1n \log u_n$.
\end{lemma}

\begin{proof}
  \leanok
  Monotonicity of the measures: $P_n(\inter S) \le P_n(S) \le P_n(\cl S)$, and $\inter S$ is a
  neighbourhood of each of its points. For the exponential bounds, $\frac1n \log(C e^{-n r}) =
  \frac{\log C}{n} - r \to -r$.
\end{proof}

\begin{lemma}[Cramér's theorem]
  \label{lem:cramer}
  \lean{LMLExplore.isLDP_cramer, LMLExplore.limsup_log_measure_le_neg_iInf_rateFunction,
    LMLExplore.neg_rateFunction_le_liminf_log_measure}
  \leanok
  \uses{def:ldp, lem:ldp_of_bounds}
  Let $P$ be a probability measure on $\R$ with a mean. The empirical means of $n$ independent
  samples of $P$ satisfy the large deviation principle with speed $n$ and the Cramér rate function
  $I(x) = \sup_t (t x - \log \E_P[e^{t X}])$ \cite[Theorem 2.2.3]{DZ98}. The lower bound
  $-I(x) \le \liminf_n \frac1n \log P_n(U)$ on the neighbourhoods $U$ of $x$ holds for every
  probability measure $P$.
\end{lemma}

\begin{proof}
  \leanok
  \uses{lem:ldp_of_bounds}
  Upper bound: LMLExplore's Cramér--Chernoff bounds on the half-lines, a closed set $F$ not
  containing the mean $m$ lying in the two half-lines bounded by $\min(F \cap [m, \infty))$ and
  $\max(F \cap (-\infty, m])$, where $I \ge \inf_F I$.
  Lower bound at a level $x$ of finite rate. If $P$ gives no mass to one side of $x$, then
  $I(x) = -\log P(\{x\})$ and the empirical mean is $x$ with probability at least $P(\{x\})^n$.
  Otherwise, as in \cite{DZ98}, truncate $P$ to $[-M, M]$: the normalized truncation $P_M$ has every
  exponential moment and charges both sides of $x$, so $t \mapsto \Lambda_M(t) - t x$ is coercive
  and its minimizer is a tilt of mean $x$; LMLExplore's exponential tilting
  (\texttt{ExponentialFamily/Cramer.lean}) gives $P_M^{\otimes n}(\abs{S_n - n x} \le n \delta) \ge
  e^{-n (I_M(x) + \varepsilon)}$ eventually, and $P^{\otimes n} \ge P([-M, M])^n P_M^{\otimes n}$.
  The functions $g_M(t) = \int_{[-M, M]} e^{t (y - x)} \, dP(y)$ increase to $g(t) = \int e^{t (y - x)}
  \, dP(y) \ge e^{-I(x)}$, and they are continuous and, for $M$ large, coercive uniformly in $M$:
  if $\inf_t g_M < e^{-I(x) - \varepsilon}$ for every $M$, the nested compact sets
  $\{t \in [-A, A] \mid g_M(t) \le e^{-I(x) - \varepsilon}\}$ have a common point $t_0$, and
  $g(t_0) = \sup_M g_M(t_0) < e^{-I(x)}$ by monotone convergence, a contradiction. Hence
  $I_M(x) - \log P([-M, M]) \le I(x) + \varepsilon$ for some $M$.
\end{proof}

\begin{lemma}
  \label{lem:rate_good}
  \lean{LMLExplore.isGoodRateFunction_rateFunction}
  \leanok
  \uses{def:good_rate}
  If the cumulant generating function of $P$ is finite on a neighbourhood of $0$, the Cramér rate
  function is a good rate function \cite[Lemma 2.2.5]{DZ98}.
\end{lemma}

\begin{proof}
  \leanok
  It is lower semicontinuous, and $I(x) \ge t \abs{x} - C$ for some $t > 0$, so its sublevel sets
  are closed and bounded.
\end{proof}

\begin{lemma}
  \label{lem:cramer_expFamily}
  \lean{LMLExplore.isLDP_cramer_expFamily}
  \leanok
  \uses{lem:cramer, lem:rate_good}
  For a nonzero measure $\rho$ on $\R$ and a mean $\mu$ of its exponential family, the empirical
  means of $n$ independent samples of the member of mean $\mu$ satisfy the large deviation
  principle with speed $n$ and the good rate function $x \mapsto \KL(x, \mu)$, extended to every
  level by the Cramér rate function of the member of mean $\mu$.
\end{lemma}

\begin{proof}
  \leanok
  \uses{lem:cramer, lem:rate_good}
  The member of mean $\mu$ has natural parameter in the interior of the natural parameters, so it
  is a probability measure whose cumulant generating function is finite near $0$.
\end{proof}

\begin{lemma}
  \label{lem:rate_continuousOn}
  \lean{MeasureTheory.Measure.continuousOn_expFamilyRate_left}
  \leanok
  For an exponential family whose statistic is not almost everywhere constant, the rate function
  $x \mapsto \KL(x, \mu)$ is continuous on the closure of the means.
\end{lemma}

\begin{proof}
  \leanok
  It is convex and lower semicontinuous, and a convex function on an interval is upper
  semicontinuous at its endpoints along the interval (LMLExplore).
\end{proof}

\begin{lemma}
  \label{lem:ldp_pi}
  \uses{def:ldp}
  If for each $k$ the measures $P^k_n$ satisfy the large deviation principle with speed $n$ and
  good rate function $I_k$, and $n_k(T) / T \to \omega_k > 0$, then the products $\bigotimes_k
  P^k_{n_k(T)}$ satisfy the large deviation principle with speed $T$ and good rate function
  $\sum_k \omega_k I_k$.
\end{lemma}

\begin{proof}
  Lower bound on the products of open sets around a point; upper bound on compact sets by a
  finite cover with products of open sets, and exponential tightness from the goodness of the
  rates.
\end{proof}

\begin{lemma}
  \label{lem:ldp_window}
  \uses{def:ldp}
  In the stack-of-rewards model, the empirical means after $N_{T,k}$ samples, with
  $\abs{N_{T,k} - T\omega_k} \le K$ and $N_{T,k}$ arbitrary otherwise, satisfy the same large
  deviation bounds as the empirical means after $\lfloor T \omega_k \rfloor$ samples.
\end{lemma}

\begin{proof}
  Upper bound: a union over the $(2K + 1)^K$ possible numbers of samples. Lower bound: the event
  that the means of the stacks after every number of samples in the window are close to the
  target has the same exponential rate.
\end{proof}

\begin{lemma}
  \label{lem:rate_ge_sq_of_subgaussian}
  \lean{ProbabilityTheory.HasSubgaussianMGF.ofReal_sq_div_le_rateFunction}
  \leanok
  If $P$ is $\sigma^2$-sub-Gaussian around its mean $m$, its Cramér rate function satisfies
  $I(x) \ge (x - m)^2 / (2 \sigma^2)$.
\end{lemma}

\begin{proof}
  \leanok
  $I(x) \ge t (x - m) - \sigma^2 t^2 / 2$, maximal at $t = (x - m) / \sigma^2$ (LMLExplore).
\end{proof}
